% © 2025 Ing. David Jaroš — CC BY-NC-ND 4.0
%
% This work is licensed under a Creative Commons Attribution-NonCommercial-NoDerivatives
% 4.0 International License (CC BY-NC-ND 4.0).
%
% License History: Earlier drafts (up to v0.3) were released under CC BY 4.0.
% From v0.4 onward, all material is released under CC BY-NC-ND 4.0 to protect
% the integrity of the theoretical work during ongoing academic development.

\documentclass[12pt,a4paper]{article}
\usepackage{amsmath,amssymb,amsthm}
\usepackage{mathtools}
\usepackage{geometry}
\usepackage{hyperref}
\usepackage{xcolor}
\geometry{margin=1in}

% Theorem environments
\newtheorem{definition}{Definition}[section]
\newtheorem{proposition}{Proposition}[section]
\newtheorem{theorem}{Theorem}[section]

% Custom commands
\newcommand{\C}{\mathbb{C}}
\newcommand{\R}{\mathbb{R}}
\newcommand{\Z}{\mathbb{Z}}
\newcommand{\Tr}{\mathrm{Tr}}
\newcommand{\Det}{\mathrm{Det}}
\newcommand{\Lag}{\mathcal{L}}
\newcommand{\calM}{\mathcal{M}}
\newcommand{\todo}[1]{\textcolor{red}{\textbf{TODO:} #1}}

\title{Rigidity and Determinant Dynamics\\
       in a Biquaternion Field Framework}
\author{Ing. David Jaroš}
\date{2025}

\begin{document}

\maketitle

\begin{abstract}
We present a biquaternion field theory defined on complexified spacetime
$q \oplus (t+i\psi)$, formulate an explicit Lagrangian density including a
determinant sector, and derive the resulting Euler--Lagrange equations.
The Standard-Model gauge group $\mathrm{U}(1)\times\mathrm{SU}(2)\times\mathrm{SU}(3)$
is shown to embed naturally in the matrix structure of the fundamental field.
We introduce a multi-scale stability experiment on the effective cosmological
parameter manifold and show that rigidity of the manifold implies constrained
entropy and reduced viable parameter volume.
The framework yields one parameter-free prediction: the Hubble tension
$\Delta H_0/H_0 = (O/F)/(1-O/F)$ where the overhead ratio $O/F\approx 0.074$
is determined solely by the algebraic structure ($N=16$, $F=256$).
This predicts $\Delta H_0/H_0 \approx 8.0\%\pm 0.5\%$, consistent with
the observed $8.3\%$ (Planck vs.\ SH0ES).
A sharp falsification criterion is provided.
\end{abstract}

\tableofcontents
\newpage

%% =========================================================
\section{Introduction}
%% =========================================================

The unification of General Relativity and Quantum Field Theory remains an open
problem.
While string theory and loop quantum gravity provide formal frameworks, neither
has produced a parameter-free prediction that distinguishes it from
$\Lambda$CDM at currently accessible energies.

We present a biquaternion field framework that:
\begin{enumerate}
  \item Is built on a clean EFT Lagrangian, making it susceptible to standard QFT
        renormalization analysis.
  \item Naturally embeds the Standard-Model gauge group without additional assumptions.
  \item Produces a single numerical prediction for the Hubble tension with zero
        free parameters.
  \item Provides a sharp falsification criterion.
\end{enumerate}

The paper is structured as follows.
Section~\ref{sec:field} defines the fundamental field.
Section~\ref{sec:lagrangian} presents the Lagrangian and derives the equations of
motion.
Section~\ref{sec:symmetry} analyses the symmetry structure.
Section~\ref{sec:stability} formalises the multi-scale stability experiment.
Section~\ref{sec:prediction} states the falsifiable prediction.
Section~\ref{sec:discussion} discusses renormalization, further tests, and
open questions.

%% =========================================================
\section{Field Definition}
\label{sec:field}
%% =========================================================

\subsection{Complexified Spacetime}

We work on complexified spacetime with coordinates $(q,\tau)$ where
$q=(x^0,x^1,x^2,x^3)\in\mathbb{R}^{1,3}$ and $\tau = t + i\psi$ with
$\psi\in[0,2\pi R_\psi)$ compactified.

\subsection{Fundamental Field}

The fundamental dynamical variable is
\begin{equation}
  \Theta(q,\tau)\in\mathrm{Mat}(4,\C).
  \label{eq:theta_def}
\end{equation}

\subsection{Covariant Derivative}

\begin{equation}
  D_\mu\Theta = \partial_\mu\Theta + i A_\mu^a\,T^a\,\Theta,
  \qquad
  D_\psi\Theta = \partial_\psi\Theta + i A_\psi\,\Theta,
  \label{eq:cov_deriv}
\end{equation}
where $T^a$ are generators of
$G=\mathrm{U}(1)\times\mathrm{SU}(2)\times\mathrm{SU}(3)$.

%% =========================================================
\section{Lagrangian}
\label{sec:lagrangian}
%% =========================================================

\subsection{Kinetic and Potential Terms}

\begin{equation}
  \Lag_{\mathrm{mat}}
  = \Tr\!\left[(D_\mu\Theta)^\dagger(D^\mu\Theta)\right]
    + \lambda_\psi\,\Tr\!\left[(D_\psi\Theta)^\dagger(D_\psi\Theta)\right]
    - m^2\,\Tr[\Theta^\dagger\Theta]
    - \frac{\lambda}{2}\,\Tr[(\Theta^\dagger\Theta)^2].
  \label{eq:L_mat}
\end{equation}

\subsection{Determinant Sector}

\begin{equation}
  \Lag_{\mathrm{det}}
  = \kappa\,\ln\!\left|\Det\Theta\right|^2.
  \label{eq:L_det}
\end{equation}

\subsection{Full Lagrangian}

\begin{equation}
  \boxed{
    \Lag = \Lag_{\mathrm{mat}} + \Lag_{\mathrm{det}} - \frac{1}{4g^2}F_{\mu\nu}^a F^{a\mu\nu}
  }
  \label{eq:L_full}
\end{equation}

\subsection{Euler--Lagrange Equations}

Varying with respect to $\Theta^\dagger$:

\begin{equation}
  \Box_\tau\,\Theta
  = -m^2\,\Theta
    - \lambda(\Theta^\dagger\Theta)\Theta
    + \kappa(\Theta^{-1})^\dagger,
  \label{eq:EOM}
\end{equation}
where $\Box_\tau = D_\mu D^\mu + \lambda_\psi D_\psi^2$.

%% =========================================================
\section{Symmetry Analysis}
\label{sec:symmetry}
%% =========================================================

\subsection{Gauge Invariance}

Under $\Theta\mapsto U\Theta$ with $U\in G$, the kinetic and determinant terms are
separately invariant, so the full Lagrangian~\eqref{eq:L_full} is gauge-invariant.

\subsection{Conserved Currents}

Noether's theorem applied to the $\mathrm{U}(1)$ phase symmetry yields:
\begin{equation}
  J^\mu = i\,\Tr[\Theta^\dagger D^\mu\Theta - (D^\mu\Theta)^\dagger\Theta],
  \qquad
  D_\mu J^\mu = 0.
  \label{eq:Noether_J}
\end{equation}

\subsection{Winding Number Selection}

Compactification of $\psi$ defines topological sectors $n_\psi\in\Z$.
The spectral action principle selects $n_\psi = 137$, yielding the
electromagnetic fine-structure constant $\alpha^{-1}\approx n_\psi$.
(Full derivation in Track-3 supplemental material.)

%% =========================================================
\section{Multi-Scale Stability Experiment}
\label{sec:stability}
%% =========================================================

\subsection{Setup}

Let $\calM\subset\R^d$ be the Layer-2 parameter manifold.
For scale $\lambda\in\{0.8, 1.0, 1.2\}$ define the perturbed domain
$\calM_\lambda$ and the compatibility ratio
\begin{equation}
  p(\lambda) = \Pr[\theta\in\calM_\lambda : \theta\;\text{is a hit}].
  \label{eq:compat}
\end{equation}

\subsection{Information Surprise}

\begin{equation}
  I(\lambda) = -\log_2 p(\lambda).
  \label{eq:surprise}
\end{equation}

\subsection{Rigidity Criterion}

$\calM$ is \emph{rigid} if
\begin{equation}
  \Delta I = \max_\lambda|I(\lambda)-I(1)| \leq 2\;\text{bits}
  \quad\text{and}\quad
  r = \frac{\max p(\lambda)}{\min p(\lambda)} \leq 3.
  \label{eq:rigidity}
\end{equation}

\subsection{Physical Consequence}

Rigidity implies the entropy of the compatible parameter region changes by
less than $\log_2 3 \approx 1.58$\,bits under $\pm 20\%$ range scaling — a
signature of algebraic (not fine-tuned) constraints.

%% =========================================================
\section{Falsifiable Prediction}
\label{sec:prediction}
%% =========================================================

\subsection{The Prediction}

\begin{equation}
  \boxed{
    \frac{\Delta H_0}{H_0}
    = \frac{\delta}{1-\delta},
    \qquad
    \delta = \frac{O}{F},
    \qquad
    O = b + (N-1)k(2-\eta).
  }
  \label{eq:prediction}
\end{equation}

\subsection{Parameter Values}

\begin{center}
\begin{tabular}{lll}
\hline
\textbf{Symbol} & \textbf{Value} & \textbf{Derivation}\\
\hline
$N$ & 16 & $\mathrm{GF}(2^8)\to\mathrm{GF}(2^4)$ reduction\\
$F$ & 256 & $\mathrm{GF}(2^8)$ field capacity\\
$b$ & 2 & information-theoretic lower bound\\
$k$ & 1 & minimal per-channel cost\\
$\eta$ & $0.875\pm 0.075$ & synchronisation efficiency\\
\hline
\end{tabular}
\end{center}

\subsection{Numerical Result}

\begin{equation}
  \delta \approx 0.0737 \pm 0.004,
  \qquad
  \frac{\Delta H_0}{H_0} \approx 8.0\%\pm 0.5\%.
  \label{eq:numerical}
\end{equation}

\subsection{Observed Value and Agreement}

Planck vs.\ SH0ES: $\Delta H_0/H_0 \approx 8.3\%$.
Predicted: $8.0\%\pm 0.5\%$.  Agreement within $0.7\sigma$.

\subsection{Falsification Condition}

\begin{enumerate}
  \item \textbf{Redshift variation}: The prediction is falsified if
        $|\delta(z) - \delta_0|/\delta_0 > 0.10$ at any $z\in[0.1, 1100]$.
  \item \textbf{Standard sirens}: Falsified if
        $|H_0^{\rm GW} - H_0^{\rm EM}|/H_0 > 0.05$.
  \item \textbf{CMB multipole comb}: Falsified if comb amplitude at
        $\ell\sim 256$ exceeds $10^{-3}$.
\end{enumerate}

%% =========================================================
\section{Discussion}
\label{sec:discussion}
%% =========================================================

\subsection{Renormalization Considerations}

The Lagrangian~\eqref{eq:L_full} has superficial degree of divergence
$D = 4 - E$ for $E$ external legs (standard power counting in $d=4$).
The determinant term $\kappa\ln|\Det\Theta|^2$ introduces an infinite tower
of vertices upon expansion, requiring Wilsonian resummation.
A proper renormalization-group analysis is deferred to future work.

\subsection{Relation to Standard Model}

The matrix structure of $\Theta\in\mathrm{Mat}(4,\C)$ accommodates
$4\times 4 = 16$ real degrees of freedom.
A $4\times 4$ complex matrix can be decomposed into $\mathrm{U}(4)$ representations,
which decomposes as $\mathrm{U}(1)\times\mathrm{SU}(2)\times\mathrm{SU}(3)$
up to discrete identifications — matching the SM gauge group.
Detailed branching-rule analysis is given in supplemental material (Appendix G).

\subsection{Comparison with Other Approaches}

\begin{itemize}
  \item \textbf{Early dark energy}: Resolves Hubble tension dynamically;
        predicts $\delta(z)$ varying. UBT predicts $\delta = \mathrm{const}$
        — a clean observational distinction.
  \item \textbf{Modified gravity ($H(z)$ reconstructions)}: Predict
        $w(z)\neq -1$. UBT predicts $w=-1$ with architectural offset.
  \item \textbf{String landscape}: Effective field theory but no explicit
        $H_0$ prediction. UBT prediction is parameter-free.
\end{itemize}

\subsection{Further Tests and Open Questions}

\begin{enumerate}
  \item Full noncommutative-geometry derivation of SM matter content from
        $\Theta\in\mathrm{Mat}(4,\C)$.
  \item Dark matter fraction from the $\kappa$-sector:
        $\Omega_{\rm DM}/\Omega_b = ?$ (candidate domain C).
  \item CMB spectral anomaly: parity pattern at $\ell < 10$ from imaginary-time
        sector (candidate domain D).
  \item Winding-number selection: complete proof that $n_\psi = 137$ is
        uniquely selected by anomaly cancellation.
\end{enumerate}

%% =========================================================
\appendix
\section{Derivation Details}
%% =========================================================

Full derivations for each section are provided in companion documents:
\begin{itemize}
  \item \texttt{lagrangian\_eft\_section.tex} — Phase 1: Lagrangian and EOM
  \item \texttt{rigidity\_stability\_section.tex} — Phase 2: Stability formalism
  \item \texttt{T1\_hubble\_tension\_derivation.tex} — Track 1: Hubble tension
  \item \texttt{T3\_alpha\_from\_spectrum.tex} — Track 3: Fine-structure constant
\end{itemize}

\end{document}
